\documentclass{article}
\usepackage{indentfirst}
\usepackage{amsmath}

\title{FLOUT -- \underline{FL}uka \underline{OU}tput processing scrip\underline{T}\newline
\newline
version 0.9.4}

\begin{document}

\maketitle

FLOUT is a shell script for processing FLUKA output binary files by means of readout codes supplied with FLUKA package.
The script requires FLUKA package distributed by CERN or INFN and bash shell version $\geq$ 4.0.
FLOUT parses FLUKA input files in FIXED and FREE formats. It also treats nested preprocessor directives.

User runs the script with the command (assume that \texttt{flout} is an alias for \texttt{flout.sh})\\
\texttt{\$ flout [mode] [file] -d [directory] [parameter 1]...[parameter n]}\\
where [mode] and [file] are mandatory arguments. All the rest arguments are optional.

Auto mode (\texttt{-a}) parses FLUKA input file,
finds all valid scoring cards and processes binary data.
Interactive mode (\texttt{-i}) parses FLUKA input file,
finds all valid scoring cards and prompts whether to process data for this card or not.
Manual mode (\texttt{-m}) processes data according to text file supplied by the user.
Scan mode (\texttt{-s}) generates text file for manual (\texttt{-m}) mode.
Optional flag (\texttt{-d}) is used to supply the path to destination directory for produced results.
The help section is displayed by running \texttt{\$ flout -h}

By default FLOUT searches for FLUKA readout programs inside CERN FLUKA
which \texttt{/bin} folder should be available on the \texttt{\$PATH}.
Supplying the \texttt{--old} parameter allows to search for readout codes inside INFN FLUKA
which requires \texttt{\$FLUPRO} shell variable.

\begin{table}[htbp]
\centering
\caption{\label{tab:tabc}Optional parameters}
\begin{tabular}{|l|l|}
\hline
Parameter             & Description                                          \\
\hline
\texttt{--log}        & save output from FLUKA readout programs to text file \\
\texttt{--clean}      & move (rather that copy) results to [directory]       \\
\texttt{--dat}        & convert USRBIN binary output to plain text           \\
\texttt{--reuse}      & reprocessing of already processed binary data        \\
\texttt{--old}        & use readout programs from old (INFN) FLUKA package   \\
\hline
\end{tabular}
\end{table}

The script has the following limitations:
\begin{enumerate}
	\item Do not use exponential notation for logical output unit in FLUKA input file.
	For instance, use 23 or 23.0 rather than 0.23E+02 or 23000.0E-03
	\item Do not \texttt{--reuse} USRYIELD data due to SIGFPE somewhere from
	\texttt{usysuw.f:924} and \texttt{usysuw.f:1410}
\end{enumerate}

\section*{Example usage}
\textbf{Task 1.} Run \texttt{example.inp}, process all binary files in auto mode.\\
\texttt{
\$ rfluka -N0 -M5 example.inp\\
\$ flout -a example.inp\\
}

\textbf{Task 2.} Run \texttt{example.inp}, process data in interactive mode,
copy it to \texttt{result} subdirectory and save log.\\
\texttt{
\$ rfluka -N0 -M5 example.inp\\
\$ mkdir result\\
\$ flout -i example.inp -d result/ --log\\
}

\textbf{Task 3.} Run \texttt{example.inp}, process unit 50 and unit 51 USRBIN data in manual mode,
convert USRBIN binary data to plain text, move it to \texttt{result} subdirectory.\\
\texttt{
\$ rfluka -N0 -M5 example.inp\\
\$ mkdir result\\
\$ cat example.txt\\
example\\
50 USRBIN\\
51 USRBIN\\
\$ flout -m example.txt -d result/ --dat --clean\\
}

\textbf{Task 4.} Make copies of \texttt{example.inp} with different seeds
and place them into subdirectories.
Run \texttt{example\_001.inp} in subdirectory \texttt{task4\_001},
run \texttt{example\_002.inp} in subdirectory \texttt{task4\_002},
process all binary files in auto mode and copy results into subdirectory \texttt{result}.
Reuse (merge) already processed data from \texttt{example\_001.inp} and \texttt{example\_002.inp}.\\
\texttt{
\$ cd task4\_001\\
\$ rfluka -N0 -M5 example\_001.inp\\
\$ flout -a example\_001.inp -d ../result\\
\$ cd ../task4\_002\\
\$ rfluka -N0 -M5 example\_002.inp\\
\$ flout -a example\_002.inp -d ../result\\
\$ cd ../result\\
\$ flout -a ../example.inp --reuse\\
}

\textbf{Task 5.} Scan \texttt{example.inp} and generate \texttt{example.txt} file for manual mode.\\
\texttt{
\$ rfluka -N0 -M5 example.inp\\
\$ flout -s example.inp\\
\$ flout -m example.txt\\
}

\section*{Changelog}
\raggedright
2022-09-12 version 0.9.4\\
-- check missing quotes, minor changes\newline

2022-09-11 version 0.9.3\\
-- add scan mode \texttt{-s} to generate input for manual mode, update docs\newline

2022-09-06 version 0.9.2\\
-- add check for writing permission in destination directory\newline

2022-09-04 version 0.9.1\\
-- add \texttt{-{}-old} parameter to support INFN FLUKA version, update docs\newline

2022-09-03 version 0.9\\
-- remove config file\newline

\end{document}
